\documentclass[preprint,12pt,authoryear]{elsarticle}

\usepackage{amsmath,amssymb}
\usepackage{booktabs}
\usepackage{graphicx}
\usepackage{xcolor}
\usepackage{listings}
\usepackage{algorithm}
\usepackage{algpseudocode}
\usepackage[hidelinks]{hyperref}

\lstset{basicstyle=\ttfamily\footnotesize, columns=fullflexible, breaklines=true, frame=single, aboveskip=6pt, belowskip=6pt}

% ------------------------------------------------------------------------------------------
% Every held-out result is generated by evaluation/report.py and evaluation/annotation.py.
% Do not type numbers into the results sections by hand.
% ------------------------------------------------------------------------------------------
% generated by evaluation/report.py - do not edit
\newcommand{\ResAnswerModel}{\texttt{--}}
\newcommand{\ResBMfiveEM}{--}
\newcommand{\ResBMfiveFOne}{--}
\newcommand{\ResBMfiveRefusal}{--}
\newcommand{\ResBMtwoEM}{--}
\newcommand{\ResBMtwoFOne}{--}
\newcommand{\ResBMtwoRefusal}{--}
\newcommand{\ResCCEM}{--}
\newcommand{\ResCCEvidenceChars}{--}
\newcommand{\ResCCFOne}{--}
\newcommand{\ResCCFinalizeRate}{--}
\newcommand{\ResCCRefusal}{--}
\newcommand{\ResCCStrictEM}{--}
\newcommand{\ResCCStrictFOne}{--}
\newcommand{\ResCCStrictRefusal}{--}
\newcommand{\ResClosedEM}{--}
\newcommand{\ResClosedFOne}{--}
\newcommand{\ResClosedRefusal}{--}
\newcommand{\ResExtractionModel}{\texttt{--}}
\newcommand{\ResFullEM}{--}
\newcommand{\ResFullFOne}{--}
\newcommand{\ResFullRefusal}{--}
\newcommand{\ResNDev}{--}
\newcommand{\ResNTest}{--}
\newcommand{\ResOracleEM}{--}
\newcommand{\ResOracleFOne}{--}
\newcommand{\ResOracleRefusal}{--}
\newcommand{\ResPropBankCommit}{\texttt{--}}
\newcommand{\ResPropBankRef}{--}
\newcommand{\ResSeed}{--}
\newcommand{\ResTriplesEM}{--}
\newcommand{\ResTriplesFOne}{--}
\newcommand{\ResTriplesRefusal}{--}

\IfFileExists{generated/annotation_macros.tex}{\input{generated/annotation_macros.tex}}{}
\providecommand{\AnnSharedParagraphs}{--}
\providecommand{\AnnFullParagraphs}{--}
\providecommand{\AnnJudgeModel}{--}

% Author notes: must all be resolved (and this macro emptied) before submission.
\newcommand{\AuthorNote}[1]{\par\noindent\fcolorbox{red}{red!5}{\parbox{0.95\linewidth}{\small\textbf{Author note:} #1}}\par}


\journal{Data \& Knowledge Engineering}

\begin{document}

\begin{frontmatter}

\title{ContextCanvas: Event-Centric Context Graphs with Semantic Roles and Context Envelopes for Grounded Multi-Hop Question Answering}

\author[inst1]{Author One\corref{cor1}}
\ead{author.one@example.org}
\cortext[cor1]{Corresponding author}
\affiliation[inst1]{organization={Affiliation},
            city={City},
            country={Country}}

\begin{abstract}
Retrieval-augmented generation passes text chunks to a language model and leaves relational structure implicit, whereas knowledge-graph approaches usually reduce text to binary triples, losing the argument structure and the temporal, spatial and provenance context of each statement.
We present ContextCanvas, a system that stores text as an event-centric property graph: each statement becomes a PropBank-grounded event frame whose arguments are typed semantic roles and which carries explicit temporal, spatial and epistemic \emph{context envelopes}.
The graph is built by a language model under a fixed frame contract and is then normalized by deterministic, independently switchable rules for kinship, place containment, aliases and role alignment.
The graph supports two answering modes: in \emph{facts mode} a language model answers from a bounded, role-labelled serialization of the subgraph around the entities named in the question; in \emph{retrieval mode} the graph ranks the source documents by their distance from those entities and the model reads the top-ranked documents as text.
On \ResNTest{} held-out MuSiQue questions, ContextCanvas reaches \ResCCEM{} exact match and \ResCCFOne{} F1, compared with \ResFullEM{} / \ResFullFOne{} for the same reader model given all candidate paragraphs and \ResBMfiveEM{} / \ResBMfiveFOne{} with BM25 retrieval.
Results are replicated with a second model family as reader, ablations quantify the contribution of role labels, context envelopes and each normalization rule; a manual annotation study measures extraction precision.
All code, prompts, configurations and per-question outputs are released.
\end{abstract}

\begin{keyword}
knowledge graphs \sep retrieval-augmented generation \sep semantic role labeling \sep PropBank \sep multi-hop question answering \sep provenance
\end{keyword}

\end{frontmatter}

\AuthorNote{Abstract numbers are filled from \texttt{generated/results\_macros.tex}. After the held-out runs, rewrite the abstract's comparative sentence to state the actual finding (better, comparable, or worse) and drop comparisons that are not significant.}

% ==========================================================================================
\section{Introduction}
\label{sec:intro}

Retrieval-augmented generation (RAG) grounds a language model's answers in retrieved text \citep{lewis2020rag}.
For questions that require combining several facts, chunk retrieval has two weaknesses.
Relations between entities remain implicit in prose, so the model must rediscover them at query time, and passages that are individually irrelevant to the question but jointly necessary are easily missed.
Graph-based variants address this by organizing extracted knowledge as a graph \citep{edge2024graphrag,gutierrez2024hipporag,guo2024lightrag}, but typically represent statements as binary triples or untyped entity links.
A triple such as \textit{(Argonne National Laboratory, employ, Ulrich Walter)} no longer records which argument is the employer, when the employment held, where it took place, or which source asserted it.

This paper argues that a knowledge representation for grounded question answering should keep three things that triples discard:
(i) the \emph{argument structure} of each statement, so that the direction of a relation is explicit;
(ii) its \emph{context} --- when and where it holds; and
(iii) its \emph{provenance and certainty}.
ContextCanvas realizes these as an event-centric property graph.
Each statement is an event node typed by a PropBank roleset \citep{palmer2005propbank}; its arguments are role-labelled edges to entity nodes; and temporal, spatial and epistemic envelopes are separate nodes attached to the event.
The graph is persisted in an embedded graph database and queried by bounded traversal from the entities named in the question.

Building such a graph with a language model is not reliable on its own.
In our development study (Section~\ref{sec:dev}) the initial prototype lost the subject of many events when a predicted role was absent from the PropBank roleset, split names such as \textit{Philippe, Duke of Orl\'eans} at the comma, and reduced \textit{son of Peter Andreas Heiberg} to \textit{Peter Andreas Heiberg}, erasing the relation.
We therefore add a deterministic normalization layer between extraction and storage, whose rules are general (they encode properties of names, dates, places and kinship expressions) and can each be disabled for ablation.

\paragraph{Contributions}
\begin{enumerate}
\item A data model for event-centric context graphs: PropBank-grounded event frames with role-labelled arguments and explicit temporal, spatial and epistemic envelopes, with a formal definition and a relational DDL for an embedded graph database (Section~\ref{sec:model}).
\item A construction pipeline in which a language model extracts events under a fixed \emph{frame contract} and a deterministic normalization layer aligns roles without discarding arguments and repairs systematic extraction errors (Section~\ref{sec:construction}).
\item A grounded answering procedure: structural anchoring of questions in the graph, hub-aware bounded evidence retrieval, a role-labelled serialization, and explicit answering policies (Section~\ref{sec:qa}).
\item An evaluation on held-out MuSiQue questions \citep{trivedi2022musique} against baselines that share the same language model and answer parser, with ablations, paired significance tests, cost accounting and an intrinsic annotation study of extraction quality (Sections~\ref{sec:setup}--\ref{sec:intrinsic}).
\end{enumerate}

% ==========================================================================================
\section{Related work}
\label{sec:related}

\paragraph{Retrieval-augmented generation}
RAG conditions generation on retrieved passages \citep{lewis2020rag}.
Long-context models can instead read all candidate passages, but their use of information depends on its position in the context \citep{liu2024lost}.
We include both a sparse-retrieval baseline \citep{robertson2009bm25} and a full-context baseline, using the same model and answer format as our system.

\paragraph{Graph-based retrieval}
GraphRAG builds an entity graph with community summaries for query-focused summarization \citep{edge2024graphrag}.
HippoRAG indexes passages through an entity graph and retrieves with personalized PageRank \citep{gutierrez2024hipporag}; LightRAG combines graph and vector retrieval \citep{guo2024lightrag}.
Think-on-Graph lets the model explore an existing knowledge graph \citep{sun2024thinkongraph}, and KAPING verbalizes retrieved triples into the prompt \citep{baek2023kaping}.
These approaches represent relations as triples or entity co-occurrence; ContextCanvas instead stores n-ary events with typed roles and context, and our ablations include a triples-only rendering of the same graph to isolate the effect of this choice.

\paragraph{Semantic representations}
Semantic role labeling assigns predicate arguments to roles \citep{gildea2002srl}; PropBank provides verb-specific rolesets \citep{palmer2005propbank}, and AMR represents full sentence meaning as rooted graphs \citep{banarescu2013amr}.
We use PropBank rolesets as a vocabulary for event types and role meanings, not as a full SRL formalism: events are produced by a language model and aligned to rolesets afterwards.

\paragraph{Event-centric and contextualized knowledge graphs}
EventKG models events and their temporal relations \citep{gottschalk2018eventkg}, and Wikidata attaches qualifiers such as validity periods and references to statements \citep{vrandecic2014wikidata}; see \citet{hogan2021kg} for a survey.
Our envelopes play the role of qualifiers, but are populated from text at ingestion time and serialized for a language model at query time.

\paragraph{Multi-hop question answering}
HotpotQA \citep{yang2018hotpotqa}, 2WikiMultiHopQA \citep{ho2020twowiki} and MuSiQue \citep{trivedi2022musique} require combining facts from several paragraphs; MuSiQue was constructed to reduce single-hop shortcuts.
Time-sensitive QA \citep{chen2021timeqa} and multi-hop knowledge editing \citep{zhong2023mquake} target temporal validity and updates, which our envelopes and discourse links are designed to support (Section~\ref{sec:limitations}).

% ==========================================================================================
\section{Data model}
\label{sec:model}

\subsection{Definition}

A \emph{context graph} is a labelled property graph
\begin{equation}
G = (V_{\mathrm{ent}} \cup V_{\mathrm{evt}} \cup V_{\mathrm{ctx}},\; E_{\mathrm{role}} \cup E_{\mathrm{ctx}} \cup E_{\mathrm{disc}}),
\end{equation}
where $V_{\mathrm{ent}}$ are entities, $V_{\mathrm{evt}}$ events and $V_{\mathrm{ctx}}$ context nodes.
An event $e \in V_{\mathrm{evt}}$ is a tuple
\begin{equation}
e = (\lambda, \sigma, \rho, \tau, \pi, \epsilon),
\end{equation}
with lemma $\lambda$, PropBank roleset $\sigma$ (e.g.\ \texttt{employ.01}), and a partial role assignment $\rho : R \rightharpoonup V_{\mathrm{ent}}$ over roles $R = \{\texttt{ARG0},\dots,\texttt{ARG5}\} \cup R_{\mathrm{mod}}$, where $R_{\mathrm{mod}}$ contains modifier roles such as \texttt{ARGM-LOC} and the spatial bridge role \texttt{location}.
Each assignment $\rho(r) = v$ is an edge $(e, r, v) \in E_{\mathrm{role}}$.
The \emph{envelopes} are optional context nodes:
a temporal envelope $\tau = (\text{expr}, t_{\text{start}}, t_{\text{end}})$ holding the raw expression and a year interval;
a spatial envelope $\pi = (\ell, \text{parent})$ holding a location and its enclosing region;
and an epistemic envelope $\epsilon = (\text{source}, c, s)$ with source identifier, confidence $c \in [0,1]$ and a speculation flag $s$.
Discourse edges relate events:
\begin{equation}
E_{\mathrm{disc}} \subseteq V_{\mathrm{evt}} \times \{\texttt{cause}, \texttt{result}, \texttt{before}, \texttt{after}, \texttt{condition}, \texttt{contrast}\} \times V_{\mathrm{evt}}.
\end{equation}

\paragraph{Identity}
Entities are identified by a \emph{matching key} $k(n)$ computed from a surface name $n$: Unicode case folding, removal of diacritics, articles and leading honorifics (only if at least two tokens remain), and trailing legal or club suffixes such as \textit{Inc.} or \textit{F.C.}.
Two mentions denote the same entity node iff their keys are equal; the first surface form seen is kept for display.
This is deliberately conservative: \textit{Adelphia} and \textit{Adelphia Communications Corporation} remain distinct.
Two events are duplicates iff they have the same roleset, the same set of (role, entity key) pairs and the same year interval; duplicates are stored once.

\paragraph{System frames}
Rolesets are taken from PropBank release \ResPropBankRef{} (commit \ResPropBankCommit{} of the frames repository).
Relations central to the evaluation domain are fixed by a registry of \emph{system frames} (Table~\ref{tab:frames}).
The extraction prompt is generated from this registry and role validation uses it with precedence over the installed PropBank release, so the roles the model is asked to produce and the roles accepted by validation cannot diverge.
Each system frame also lists the verbs that may carry it (for \texttt{locate.01}: \textit{locate, be located, situate, lie, headquarter, take place}, \ldots).
The contract applies in both directions: a listed verb whose extracted sense does not exist in PropBank (\textit{co-found} as \texttt{co-found.01}) takes its system frame, and an extracted event whose verb is not listed keeps its arguments but receives the verb's own sense instead: language models attach system frames to verbs that do not express them (\textit{hold} or \textit{be known as} as \texttt{locate.01}, \textit{be home to} as \texttt{member.01}), and because containment is transitive, a single false \texttt{locate.01} fact corrupts every place chain through it.
Two internal frames are produced only by the normalization layer: \texttt{alias.01} (two names of one entity) and \texttt{affiliate.01} (possessive association).

\begin{table}[t]
\centering\small
\caption{System frames. The extraction prompt is generated from this table, and role validation uses it with precedence over PropBank.}
\label{tab:frames}
\begin{tabular}{lll}
\toprule
Frame & \texttt{ARG0} & \texttt{ARG1} / \texttt{ARG2} \\
\midrule
\texttt{locate.01} & -- & thing located / location \\
\texttt{own.01} & owner & possession \\
\texttt{parent.01} & parent & child \\
\texttt{marry.01} & spouse & other spouse \\
\texttt{sibling.01} & person & sibling \\
\texttt{bear.02} & parent giving birth & person born \\
\texttt{employ.01} & employer & employee / position \\
\texttt{member.01} & member & group \\
\texttt{found.01} & founder & organisation founded \\
\texttt{distribute.01} & distributor & thing distributed / recipient \\
\texttt{manufacture.01} & manufacturer & product \\
\texttt{partner.01} & partner & other partner (not marriage) \\
\midrule
\texttt{alias.01}$^*$ & entity & alternative name \\
\texttt{affiliate.01}$^*$ & possessor & associated entity \\
\bottomrule
\end{tabular}
\par\smallskip{\footnotesize $^*$Internal frames, produced only by the normalization layer.}
\end{table}

\subsection{Storage}

The graph is stored in K\`uzu, an embedded property-graph database \citep{feng2023kuzu}, with the schema of Listing~\ref{lst:ddl}.
Each event, its role edges and its envelopes are written in a single transaction.
For traversal, an in-memory mirror of the graph is maintained and rebuilt from the database when a database is opened, so a persisted graph remains queryable across sessions; entity keys and event signatures are indexed so that resolution and de-duplication take constant expected time per mention.

\begin{lstlisting}[caption={Schema (K\`uzu DDL).},label={lst:ddl}]
CREATE NODE TABLE Entity(id STRING, name STRING, PRIMARY KEY (id));
CREATE NODE TABLE Event(id STRING, lemma STRING, sense_id STRING, PRIMARY KEY (id));
CREATE NODE TABLE TimeContext(id STRING, raw_expr STRING, start_year INT64,
                              end_year INT64, PRIMARY KEY (id));
CREATE NODE TABLE PlaceContext(id STRING, location_name STRING,
                               parent_region STRING, PRIMARY KEY (id));
CREATE NODE TABLE EpistemicContext(id STRING, source STRING, confidence DOUBLE,
                                   is_speculative BOOLEAN, PRIMARY KEY (id));
CREATE REL TABLE ArgumentRole(FROM Event TO Entity, role STRING);
CREATE REL TABLE InTemporalContext(FROM Event TO TimeContext);
CREATE REL TABLE InSpatialContext(FROM Event TO PlaceContext);
CREATE REL TABLE InEpistemicContext(FROM Event TO EpistemicContext);
CREATE REL TABLE DiscourseLink(FROM Event TO Event, relation STRING);
\end{lstlisting}

% ==========================================================================================
\section{Graph construction}
\label{sec:construction}

\subsection{Extraction}

Paragraphs longer than 900 characters are split at sentence boundaries, and the document title is prefixed to later chunks to support pronoun resolution.
For each chunk a language model emits JSON listing up to 15 events, each with a lemma, a roleset identifier, role-to-name assignments, optional temporal and spatial context, and optional discourse links.
The prompt (released with the code) lists the system frames of Table~\ref{tab:frames}, instructs the model to use them only for the relations they denote and to express those relations with these verbs rather than with a copula and a relational noun, and asks for entity names rather than phrases containing them (\textit{Acme Corp} rather than \textit{Acme Corp's support}).
The source identifier is not shown to the model, so extraction is a deterministic function of the text at temperature~0; outputs are cached by a hash of the model, prompt and text, which lets all system variants in our experiments share one set of extractions.
A schema layer coerces common deviations (alternative role names, list-valued roles, a sense identifier written into the lemma field, truncated JSON) into the data model.

\subsection{Normalization}

Each extracted event passes through the following rules before storage.
Each rule can be disabled independently (Section~\ref{sec:ablation}).

\begin{description}
\item[Frame--verb contract.] A system frame on a verb not listed for it is replaced by the verb's own sense, and a listed verb with a non-existent sense takes its frame (Section~\ref{sec:model}).
\item[Date migration.] Arguments that are pure temporal expressions (recognized by a vocabulary of month names, years, era markers and temporal function words) move into the temporal envelope, so dates never become entities.
\item[Placeholder filtering.] Arguments such as \textit{unknown}, \textit{someone} or bare pronouns are removed, as are participants the extractor itself marks as inferred (\textit{network executives (implied)}); otherwise a single \textit{unknown} node links unrelated events.
\item[Self-relations.] No entity relates to itself. A kinship event between a person and themselves is discarded; in any other event, an entity that fills two core roles keeps only the non-agent one, so that \textit{DeSoto Records founded DeSoto Records in 1989} is stored as \textit{DeSoto Records was founded in 1989}.
\item[Kinship.] Kinship is rewritten into the direction-explicit frames \texttt{parent.01}, \texttt{marry.01} and \texttt{sibling.01}. Kinship lemmas are grouped by what the first argument is relative to the others (child-of, parent-of, spouse, sibling), and relational noun phrases in arguments (\textit{elder sister of X}, \textit{son of X and Y}, or \textit{younger son} with a separate argument \textit{X}) are parsed into the same frames. Plural phrases (\textit{children of \ldots}) name groups and are not rewritten, and phrases inside birth events are not rewritten.
\item[Copula and birth.] Subject--predicate copulas are mapped onto \texttt{be.01} (\texttt{ARG1} topic, \texttt{ARG2} comment). Birth events are mapped onto \texttt{bear.02} with the person born as \texttt{ARG1}; one remaining argument is the birthplace, two or more are parents.
\item[Place hierarchies.] A comma-qualified place in a locative slot (\textit{Canyon, Texas, United States}) becomes its head plus explicit containment events (\textit{Canyon} in \textit{Texas}, \textit{Texas} in \textit{United States}). Strings containing digits, lower-case fragments, title words (\textit{Duke}, \textit{Prince}) or more than five parts are left intact, so names such as \textit{Philippe, Duke of Orl\'eans} and quantities such as \textit{161,000 customers} are never split.
\item[Aliases.] Definitions of the form \textit{Long Name (ABBR)} in the source text create \texttt{alias.01} events if the abbreviation is alphabetic and spells the initials of the name; the name is trimmed to the shortest suffix that does, and never extends across a sentence boundary. Unaligned abbreviations (e.g.\ in another language) are not mined: a wrong alias merges unrelated entities, while a missing one only loses a link.
\item[Possessive links.] \textit{X's Y}, where both parts are proper names and $X$ is an entity in its own right, yields \texttt{affiliate.01}$(X, \textit{X's Y})$. Titles containing a possessive (\textit{Grant's First Stand}) are excluded. When the possessed part is a common noun phrase (\textit{Orion Pictures' support}), the argument is kept as stated and $X$ is attached to the same event as a possessor role; replacing the argument by $X$ would change the fact (\textit{ended Algeria's status as a colony} is not \textit{ended Algeria}).
\end{description}

\paragraph{Sense selection and role alignment}
If the predicted roleset is unknown, a sense is chosen among the rolesets of the lemma by word overlap between the paragraph and the roleset's name and role descriptions.
Roles are then aligned to the roleset (Algorithm~\ref{alg:align}).
Discarding arguments that a roleset does not declare --- the conventional validation --- disconnects their entities from the graph: in our development data it removed the subject of every \texttt{be}, \texttt{consist} and \texttt{star} event.
Our alignment instead maps the arguments onto the declared slots, never fills an empty \texttt{ARG0} from another position (agency cannot be inferred from argument position: a passive \textit{owned by X} and a location \textit{in X} have the same shape), and drops only high-numbered arguments that fit nowhere.

\begin{algorithm}[t]
\caption{Role alignment of proposed core roles $P$ to the declared core roles $D$ of a roleset.}
\label{alg:align}
\begin{algorithmic}[1]
\If{$D = \emptyset$} \Return $P$ \EndIf
\If{$\texttt{ARG0} \in P$ and $\texttt{ARG0} \notin D$} \Comment{unaccusative or copular frame}
  \State map the $i$-th smallest role of $P$ to the $i$-th smallest role of $D$
\Else
  \State keep every $r \in P \cap D$
  \State $F \gets$ declared roles not yet used, excluding \texttt{ARG0}
  \For{each $r \in P \setminus D$ in ascending order}
    \If{$F \neq \emptyset$} move $r$ to the smallest role in $F$ and remove it from $F$
    \ElsIf{$r \in \{\texttt{ARG0}, \texttt{ARG1}, \texttt{ARG2}\}$} keep $r$ unchanged
    \Else\ drop $r$
    \EndIf
  \EndFor
\EndIf
\State keep all modifier and context roles
\end{algorithmic}
\end{algorithm}

% ==========================================================================================
\section{Question answering}
\label{sec:qa}

\subsection{Anchoring}

The entities named in the question are located by matching graph entity names against the question, after normalizing possessives on both sides, in four tiers:
(1) word-boundary occurrences of an entity name;
(2) a capitalized question span that is a word-boundary prefix of an entity name (\textit{Caroline LeRoy} $\to$ \textit{Caroline LeRoy Webster}; \textit{The Unwinding} $\to$ \textit{The Unwinding: An Inner History \ldots});
(3) a capitalized span whose words form an ordered subsequence of an entity name with the same first and last word and at most two extra words (\textit{Frances Tupper} $\to$ \textit{Frances Am\'elia Tupper});
and, only if these fail, (4) an entity whose name contains all words of a multi-word capitalized span.
Overlapping matches are resolved in favour of capitalized, longer and exact matches, and lower-case matches of generic nouns (\textit{employer}, \textit{author}) are used only if nothing name-like matches.
If several anchors remain, anchors attached to more than $h$ events (\emph{hubs}, $h = 25$) are dropped, so that \textit{Serbia} in \textit{Tihomir of Serbia} does not flood the evidence.

\subsection{Evidence retrieval and serialization}

Evidence is collected by breadth-first search from the anchors, where one hop is an entity--event--entity step and the search is limited to three hops.
Hub entities other than the anchors are reached but not expanded, and \texttt{alias.01} edges are crossed without counting a hop.
Events are ordered by distance from the anchors and, within a distance, by word overlap between the question and the event's lemma and role descriptions, and are serialized until a budget derived from the model's context window is reached.
Each event is rendered with its role meanings and envelopes, for example:

\begin{lstlisting}
- [ev4_60] employ (employ.01): ARG0[employer]=German Aerospace Center;
    ARG1[employee]=Ulrich Walter; ARG2[position]=trainee
  Temporal Scope: 1988 to 1990 [1988-1990]
\end{lstlisting}

The serialization is followed by up to ten \emph{path candidates}: entities reached by the search, selected in turn across distances and listed with the role-labelled path from the anchor.
Epistemic envelopes are shown only when informative (confidence below one or speculative); source identifiers are always retained internally and used to compute supporting-paragraph predictions.

\subsection{Graph-guided retrieval}
\label{sec:retrieval}

Extraction loses information: a relation the extractor missed cannot be read from the facts.
In retrieval mode the graph is used only to select text.
The anchors are found as above, and each source document is linked to the graph search in three ways, of which it takes the closest: through an event extracted from it (at the event's distance); through its title, when the title names an entity the search reached (\textit{Pecos County, Texas} names \textit{Pecos County}; a parenthetical qualifier is ignored); and through a \emph{mention}, when its text contains the name of an entity the search reached.
Mention links compensate for extraction misses, and they matter for multi-hop questions in particular: the document that answers the second hop typically mentions the bridge entity without being about it, and nothing about that entity may have been extracted from it.
Only names of at least four characters containing a capital letter create mention links, and hub entities other than the anchors create none, as in the graph search.
At equal distance, event links rank before title links and title links before mention links.
The $k$ nearest documents (ties in ingestion order) are given to the reader as text, with the same prompt as the text baselines.
If the graph reaches fewer than $k$ documents, the remaining slots are filled by BM25, so that the evidence budget equals that of the BM25 top-$k$ baseline; a variant without filling isolates the graph's contribution.
Retrieval mode keeps the provenance of facts mode: every selected document is justified by a path from the question's entities.

\subsection{Answering}

The language model receives the serialized evidence and the question and must end its output with one of
\begin{center}
\texttt{FINAL ANSWER: <entity>}\qquad or\qquad \texttt{FINAL ANSWER: NOT\_IN\_EVIDENCE}.
\end{center}
The refusal criterion is stated once and depends on the \emph{answering policy}:
under the \emph{strict} policy the model refuses unless the facts state every relation the question asks about;
under the \emph{best-supported} policy it identifies the requested answer type and follows the chain of facts to the best-supported entity of that type, refusing only if no chain reaches one.
If a generation ends without a final line (typically by reaching the token limit), one short follow-up call asks for the final line given the reasoning already produced, so that truncation is not counted as refusal.
Answers that resolve to an evidence entity are returned in that entity's exact surface form.

% ==========================================================================================
\section{Experimental setup}
\label{sec:setup}

\paragraph{Data}
We use the answerable development set of MuSiQue \citep{trivedi2022musique}.
Each question comes with 20 paragraphs, of which some support the answer and the rest are distractors.
For every question a fresh graph is built from its 20 paragraphs; this matches the benchmark's distractor setting.
The first \ResNDev{} questions of the file were used during development (Section~\ref{sec:dev}) and are excluded.
From the remaining questions, \ResNTest{} were sampled once with seed \ResSeed{} and fixed in a manifest; every system is evaluated on exactly these questions.

\paragraph{Models}
ContextCanvas uses two models with different roles.
The \emph{extractor} (\ResExtractionModel{}) builds the graph; extraction is a high-volume structured task with one call per paragraph chunk.
The \emph{reader} (\ResAnswerModel{}) answers questions; it is also the only model used by the text baselines, so every system answers with the same model and the comparison isolates the representation of the evidence.
All models are served locally through Ollama at temperature~0 with a context window of 8{,}192 tokens, of which 2{,}048 are reserved for the answer.
Section~\ref{sec:models} replaces the reader with a model of another family, replaces the extractor with a smaller model of the same family, and uses a single small model for both roles.
The prompts and the full configuration of every run are released with the per-question outputs.
\AuthorNote{Cite the technical reports of the models used.}

\paragraph{Systems}
All systems share the answer format, the answer parser and the evidence budget.
\emph{Closed-book}: the question alone.
\emph{BM25 top-$k$}: the $k \in \{2, 5\}$ paragraphs ranked highest by BM25 \citep{robertson2009bm25}.
\emph{Full-context}: all 20 paragraphs, truncated to the evidence budget.
\emph{Oracle}: only the gold supporting paragraphs; this is an upper bound for text-based reading, not a competitor.
\emph{ContextCanvas} in facts mode under the best-supported and the strict policy, and in retrieval mode with $k=5$: with BM25 filling, without it, and without mention links.

\paragraph{Metrics}
Exact match (EM) and token F1 after SQuAD normalization \citep{rajpurkar2016squad}, maximized over the gold answer and its aliases; supporting-paragraph F1, where a system's predicted support is the set of paragraphs whose content was shown to the answerer (for ContextCanvas: the sources of the serialized events); refusal rate; and cost, split into index-time tokens (extraction) and query-time tokens (answering) per question, with wall-clock model time.

\paragraph{Statistics}
We report 95\% percentile-bootstrap confidence intervals \citep{efron1993bootstrap}.
Systems are compared on the same questions with an exact McNemar test on EM \citep{mcnemar1947} and a paired bootstrap test on F1 \citep{dror2018hitchhiker}.
Ablation $p$-values are adjusted for multiple comparisons with the Holm procedure \citep{holm1979}.

% ==========================================================================================
\section{Results}
\label{sec:results}

\subsection{Main comparison}

\begin{table}[t]
\centering\small
\caption{Held-out MuSiQue results. All systems use the same language model, answer format and evidence budget.}
\label{tab:main}
\resizebox{\ifdim\width>\linewidth\linewidth\else\width\fi}{!}{% generated by evaluation/report.py - do not edit
\begin{tabular}{lrrrrr}
\toprule
System & EM [95\% CI] & F1 [95\% CI] & Sup.\ F1 & Refusal \% & $p$ (EM) \\
\midrule
\multicolumn{6}{c}{(no results yet)} \\
\bottomrule
\end{tabular}
}
% generated by evaluation/report.py - do not edit
\par\smallskip{\footnotesize $n=0$ held-out questions. CIs: percentile bootstrap (10{,}000 resamples). $p$: exact McNemar test on per-question EM against ContextCanvas (best-supported). $^\dagger$Upper bound: receives the gold supporting paragraphs.}

\end{table}

\begin{table}[t]
\centering\small
\caption{Cost per question.}
\label{tab:cost}
% generated by evaluation/report.py - do not edit
\begin{tabular}{lrrrr}
\toprule
System & Index tok./q & Query tok./q & Index s/q & Query s/q \\
\midrule
\multicolumn{5}{c}{(no results yet)} \\
\bottomrule
\end{tabular}

% generated by evaluation/report.py - do not edit
\par\smallskip{\footnotesize Mean language-model tokens and seconds per question. Index: extraction of the question's paragraphs (cached extractions are charged their original cost). Query: answering, including follow-up calls for truncated answers.}

\end{table}

Table~\ref{tab:main} reports the main comparison.
ContextCanvas (best-supported) obtains \ResCCEM{} EM and \ResCCFOne{} F1; the strict policy obtains \ResCCStrictEM{} EM with a refusal rate of \ResCCStrictRefusal\%, against \ResCCRefusal\% for the best-supported policy.
The full-context baseline obtains \ResFullEM{} EM, BM25 top-5 \ResBMfiveEM{}, the closed-book model \ResClosedEM{}, and the oracle \ResOracleEM{}.
On average the serialized evidence comprised \ResCCEvidenceChars{} characters, and the follow-up call for truncated answers was needed for \ResCCFinalizeRate\% of questions.
Table~\ref{tab:cost} separates index-time from query-time cost: ContextCanvas pays for extraction once per paragraph, whereas the text baselines read paragraphs at every query.

\AuthorNote{Decide which mode is the headline system and regenerate the tables with \texttt{report.py --reference} set to it. Interpret Table~\ref{tab:main} once it is filled: (o) facts mode versus retrieval mode, and retrieval mode versus BM25 at the same budget; (a) whether ContextCanvas differs significantly from full-context and BM25; (b) the precision/coverage trade-off between the two policies; (c) the gap to the oracle; (d) the cost profile of Table~\ref{tab:cost}. Do not claim superiority where $p \geq 0.05$.}

\subsection{Effect of reasoning depth}

\begin{table}[t]
\centering\small
\caption{Exact match by number of hops.}
\label{tab:hops}
% generated by evaluation/report.py - do not edit
\begin{tabular}{l}
\toprule
System \\
\midrule
\multicolumn{1}{c}{(no results yet)} \\
\bottomrule
\end{tabular}

% generated by evaluation/report.py - do not edit
\par\smallskip{\footnotesize Exact match (\%) by number of reasoning hops.}

\end{table}

Table~\ref{tab:hops} breaks EM down by the number of composed single-hop questions.
\AuthorNote{If the sample contains few 3- and 4-hop questions, report their counts and treat the trend as indicative; consider stratifying the sample by hop count.}

\subsection{Robustness to the language model}
\label{sec:models}

\begin{table}[t]
\centering\small
\caption{Headline systems under different model configurations, on the same held-out questions.}
\label{tab:models}
\resizebox{\ifdim\width>\linewidth\linewidth\else\width\fi}{!}{% generated by evaluation/report.py - do not edit
\begin{tabular}{lllrr}
\toprule
Configuration & Extractor & Reader & CC EM & Full EM \\
\midrule
\multicolumn{5}{c}{(no results yet)} \\
\bottomrule
\end{tabular}
}
\IfFileExists{generated/tab_models_note.tex}{% generated by evaluation/report.py - do not edit
}{}
\end{table}

Table~\ref{tab:models} repeats the main comparison with a reader from another model family, with a smaller extractor, and with a single small model in both roles.
Configurations that share the extractor share the extracted graphs exactly (extractions are cached by extractor, prompt and text), so differences between them are due to the reader alone.
\AuthorNote{State whether the ranking of ContextCanvas against the full-context baseline is stable across readers, and what the small configuration shows about the model size the method needs.}

\subsection{Ablations}
\label{sec:ablation}

\begin{table}[t]
\centering\small
\caption{Ablations. Each variant removes one component from the full system.}
\label{tab:ablation}
% generated by evaluation/report.py - do not edit
\begin{tabular}{lrrrrrr}
\toprule
Variant & EM & $\Delta$EM & $p_{\text{Holm}}$ & F1 & $\Delta$F1 & $p_{\text{Holm}}$ \\
\midrule
\multicolumn{7}{c}{(no results yet)} \\
\bottomrule
\end{tabular}

% generated by evaluation/report.py - do not edit
\par\smallskip{\footnotesize Each row removes one component from the full system; all variants share the same cached extractions. EM: exact McNemar; F1: paired bootstrap; both Holm-corrected over all ablations.}

\end{table}

Table~\ref{tab:ablation} removes one component at a time.
All variants reuse the same cached extractions, so differences are attributable to normalization, retrieval or serialization alone.
The \emph{triples only} variant serializes the same graph as subject--predicate--object triples without role meanings or envelopes, and path candidates without role labels; it isolates the contribution of the event-centric representation from that of the extraction and anchoring machinery.
\AuthorNote{Discuss which components matter. Normalization rules that fire rarely may show small, non-significant average effects; report how often each rule fired (counts can be computed from the stored graphs) alongside its effect.}

\subsection{Error analysis}

\begin{table}[t]
\centering\small
\caption{Outcome attribution for ContextCanvas (best-supported) on the held-out questions.}
\label{tab:failures}
% generated by evaluation/report.py - do not edit
\begin{tabular}{lrr}
\toprule
Outcome & Questions & \% \\
\midrule
\multicolumn{3}{c}{(no results yet)} \\
\bottomrule
\end{tabular}

% generated by evaluation/report.py - do not edit
\par\smallskip{\footnotesize Outcome attribution for ContextCanvas. `answer in evidence' means a gold alias occurs in the evidence shown to the answerer.}

\end{table}

Each question is attributed to the stage that determined its outcome (Table~\ref{tab:failures}).
A refusal or wrong answer is distinguished by whether a gold answer string occurred in the evidence shown to the model: if it did not, the error lies in extraction or retrieval; if it did, in anchoring, ranking or reasoning.
\AuthorNote{Summarize the dominant categories and give two or three concrete examples per category from \texttt{predictions.jsonl}.}

% ==========================================================================================
\section{Development study}
\label{sec:dev}

The system was developed on the first 50 answerable questions of the MuSiQue development file; these questions are excluded from all results above, and the numbers in this section are \emph{not} evidence of generalization.
We report them because the development process exposed failure modes that are likely to affect any LLM-to-graph pipeline.

The first complete prototype answered 50.0\% of the development questions exactly (53.9 F1).
Analysis of its traces identified general defects, each fixed by a general rule rather than by question-specific changes:
role validation that discarded arguments undeclared in the roleset and thereby disconnected subjects;
comma splitting of all arguments, which broke names and numbers;
name normalization that removed relational nouns (\textit{son of X} $\to$ \textit{X}) and so erased kinship;
anchoring on generic nouns and on single tokens of multi-word names;
unbounded evidence that traversed hubs such as compass directions;
and answer parsing that treated truncated generations as refusals.
After these changes the system answered 68.0\% of the development questions exactly (75.0 F1).
Of the 16 remaining development failures, six lacked the answer relation in every extracted fact, two contained extraction errors that made the evidence wrong, and four involved gold answers that were ambiguous or differed in surface form from an equally correct prediction.
The remaining four (three generations truncated before the final line and one question anchored on a hub entity) motivated the final truncation and hub-pruning rules; their effect was not re-measured on the development questions.

% ==========================================================================================
\section{Extraction quality}
\label{sec:intrinsic}

\begin{table}[t]
\centering\small
\caption{Manual and LLM-judged assessment of extracted events.}
\label{tab:annotation}
\IfFileExists{generated/tab_annotation.tex}{\resizebox{\ifdim\width>\linewidth\linewidth\else\width\fi}{!}{\input{generated/tab_annotation.tex}}\input{generated/tab_annotation_note.tex}}{(annotation pending)}
\end{table}

End-to-end accuracy conflates extraction, retrieval and reasoning.
To measure extraction directly, we sampled \AnnFullParagraphs{} paragraphs from the extraction cache of the held-out questions.
For every extracted event, annotators judged whether the paragraph states the relation between the given arguments, whether every argument is in the correct role, and whether an attached time or place is correct; for every paragraph they counted relations between named entities that were stated but not extracted.

The study combines human and model annotation.
Two human annotators, working independently after a joint calibration on development paragraphs, annotated a random subset of \AnnSharedParagraphs{} paragraphs.
An LLM judge (\AnnJudgeModel{}), from a different model family than the extractor and the readers, annotated the full sample.
Its instructions are the annotation guidelines given to the humans, verbatim; it received the same information, including the meaning of each role; and its outputs were checked against the structure the guidelines require, with conditional judgments removed where the guidelines leave them blank.
On the shared subset we report Cohen's $\kappa$ \citep{cohen1960kappa} between the two humans and between each human and the judge.
The judge's measures on the full sample are interpreted only to the extent that its agreement with the humans is comparable to the agreement between the humans.
Table~\ref{tab:annotation} reports event precision, role-direction accuracy, time/place accuracy and missed facts per paragraph (a proxy for recall).
The guidelines, the judge's prompt and settings, and every raw judge response are released.
\AuthorNote{Describe the human annotators (background, involvement in the project) and the time spent. State whether judge--human agreement approaches human--human agreement for each measure; if it does not for missed facts, report recall from the humans only. Note any systematic leniency of the judge relative to the humans.}

% ==========================================================================================
\section{Storage and software invariants}
\label{sec:storage}

\begin{table}[t]
\centering\small
\caption{Storage behaviour on synthetic employment events.}
\label{tab:storage}
% generated by evaluation/report.py - do not edit
\begin{tabular}{rrrrr}
\toprule
Events & Inserts/s & Disk bytes & Bytes/event & Median query ms \\
\midrule
\multicolumn{5}{c}{(no results yet)} \\
\bottomrule
\end{tabular}

% generated by evaluation/report.py - do not edit
\par\smallskip{\footnotesize Synthetic employment events; queries are served from the in-memory traversal mirror.}

\end{table}

Table~\ref{tab:storage} reports insertion throughput, on-disk size and query latency for synthetic graphs of increasing size.
Queries are served from the in-memory mirror, so latency reflects traversal rather than database access.
The released test suite checks invariants of the implementation --- transactional insertion, envelope serialization, discourse links, role alignment, the normalization rules and the evaluation statistics --- on hand-constructed inputs.
These tests verify that components behave as specified; they are not measurements of reasoning accuracy.

% ==========================================================================================
\section{Limitations and threats to validity}
\label{sec:limitations}

\paragraph{Single benchmark, few models}
Results are for one English Wikipedia-based benchmark and a small number of locally served open-weight models; larger or proprietary readers may narrow or widen the differences between systems.
The temporal and epistemic envelopes are largely unexercised by MuSiQue, whose questions rarely depend on time or source reliability; claims about their usefulness for temporal reasoning or knowledge updates require benchmarks such as time-sensitive QA \citep{chen2021timeqa} or multi-hop knowledge editing \citep{zhong2023mquake}, which we leave to future work.

\paragraph{Per-question graphs}
Following the distractor setting, each question has its own small graph.
Building one graph over a large corpus would stress entity resolution, whose conservative matching keys would then leave more co-referent names unmerged, and would change the cost profile.

\paragraph{Development on the same benchmark}
The normalization rules were designed while inspecting failures on 50 MuSiQue development questions.
Although the rules encode general properties of names, places and kinship expressions and the evaluation questions are disjoint from the development questions, the rules may favour the entity types frequent in this benchmark (people, places, organisations, creative works).

\paragraph{Answer format}
Exact match penalizes correct answers in a different surface form (\textit{Metropolitan Borough of Knowsley} for \textit{Knowsley}).
We report F1 alongside EM and do not apply benchmark-specific answer rewriting beyond removing a trailing region qualifier when the question does not ask for a region.

\paragraph{Heuristics and language}
Normalization relies on English morphology and capitalization.
Sense selection uses word overlap and is a weak disambiguator; the system frames therefore fix the senses of the relations most important for the evaluation.

\paragraph{Cost}
Extraction requires one model call per paragraph chunk, a cost paid at indexing time and amortized over all questions that use the graph.

% ==========================================================================================
\section{Conclusion}
\label{sec:conclusion}

ContextCanvas represents text as an event-centric property graph in which statements keep their argument structure, context and provenance, and answers questions from a bounded, role-labelled serialization of that graph.
Making this work with a language model as extractor required a deterministic normalization layer whose central design choice is to never discard an argument it cannot place: an imprecise role label is recoverable at query time, while a dropped entity is not.
\AuthorNote{State the main empirical finding in one or two sentences once the held-out results are in.}

% ==========================================================================================
\section*{Data and code availability}
Code, prompts, configuration files, the evaluation manifest (question identifiers and seed), per-question predictions with model outputs, and the annotation files are available at \url{https://example.org/contextcanvas}.
PropBank frames: release \ResPropBankRef{}, commit \ResPropBankCommit{}, recorded with its archive checksum in the released lock file.
\AuthorNote{Replace with the archived repository URL (e.g.\ a Zenodo DOI).}

\section*{Declaration of competing interest}
The authors declare no competing interests.
\AuthorNote{Confirm or replace.}

\section*{Declaration of generative AI use}
\AuthorNote{The LLM judge is part of the research method and is described in Section~\ref{sec:intrinsic}; this declaration concerns the writing process. Elsevier requires a statement if generative AI tools were used in preparing the manuscript. Describe truthfully which tools were used, for which parts (e.g.\ code development, drafting of text), and confirm that the authors reviewed and edited all content and take responsibility for it.}

\bibliographystyle{elsarticle-harv}
\bibliography{references}

\end{document}
